% Generated draft: requires mathematical and source review.
\documentclass[11pt]{article}
\usepackage[margin=25mm]{geometry}
\usepackage{fontspec}
\setmainfont{Latin Modern Roman}
\usepackage{amsmath,amssymb,mathtools,unicode-math}
\setmathfont{Latin Modern Math}
\usepackage{enumitem,needspace,xurl,hyperref}
\setlength{\parindent}{0pt}
\setlength{\parskip}{5pt}
\setlength{\emergencystretch}{3em}
\newcommand{\N}{\mathbb{N}}
\newcommand{\Z}{\mathbb{Z}}
\newcommand{\Q}{\mathbb{Q}}
\newcommand{\R}{\mathbb{R}}
\newcommand{\C}{\mathbb{C}}
\newcommand{\F}{\mathbb{F}}
\newcommand{\A}{\mathbb{A}}
\newcommand{\PP}{\mathbb{P}}
\newcommand{\E}{\mathbb{E}}
\newcommand{\eps}{\varepsilon}
\newcommand{\dd}{\,\mathrm{d}}
\begin{document}
\section*{UnsolvedMath ID: 2307027}
\textbf{Source problem code:} catalogue entry
\subsection{Definitions and mathematical statement}
Let $D\subset\C$ be a domain, meaning a nonempty connected open set. For fixed $x,y\in D$, define
\[
\rho_D(x,y)=\sup\{|f(x)-f(y)|:
 f\text{ holomorphic on }D,\quad |f'(z)|\leq1\text{ for every }z\in D\}.
\]
Let the intrinsic Euclidean path distance be
\[
L_D(x,y)=\inf\{\operatorname{length}(\gamma):
\gamma\text{ is a rectifiable path in }D\text{ joining }x\text{ to }y\}.
\]
The source writes paths in $G$, although only $D$ is introduced; the display makes the natural $G=D$ reading explicit.

The immediate comparison is
\[
|x-y|\leq\rho_D(x,y)\leq L_D(x,y).
\]
The lower bound uses the admissible function $f(z)=z$; the upper bound is the source's path-length comparison. If $D$ is convex, straight segments give $\rho_D(x,y)=|x-y|$ for every pair.

L. A. Rubel asks for a geometric description or estimates of $\rho_D$ for general domains. The question does not state one proposed formula, comparison constant, or domain regularity class. The source also says the Euclidean equality holds ``but not otherwise.'' Its precise intended quantifier and domain qualifications are absent; a necessity-of-convexity reading needs special care with removable boundary sets and is not assumed here. No simple-connectedness or boundary regularity restriction is silently added.

The supplied summary reports the question open as of the 2018 edition and no definitive resolution located. The target remains to relate the holomorphic derivative-constrained distance to the actual geometry and analytic removability properties of the domain, with the desired strength of description not specified by the source.
\subsection{Short English statement}
Rubel asks how the largest separation attainable by holomorphic functions with derivative bounded by one depends on the domain. It lies between Euclidean and shortest-path distances and equals Euclidean distance on convex domains. The source gives no exact general conjecture, and its converse-sounding convexity remark needs qualification.
\subsection{Sources}
\begin{itemize}
\item Local UnsolvedMath registry, permanent ID 2307027, exact source code catalogue entry, ``Research Problems in Function Theory --- Problem 7.27''; status open, classification OPEN-TRIAGE.
\item Hayman, \emph{Research Problems in Function Theory} (2018), Problem 7.27, attributed to L. A. Rubel, \url{https://arxiv.org/abs/1809.07200}. The record identifies extraction from public arXiv source TeX, accessed July 29, 2026; the 2018 update was retained for status triage.
\item The supplied summary reports openness as of the 2018 edition and no definitive resolution located. No additional authoritative citation, exact geometric characterization, or precise search date is supplied.
\end{itemize}
\end{document}
